\section{Conclusion}

Digit-distribution tests can be part of an election-forensics workbench, but
only in a narrow and carefully labeled role. This paper defines that role for
RCOUNT packages: extract stated digit features from hash-bound summaries or
large count tables, compare them to declared references, report a statistic or
rank, and attach the weak-signal caveat required by the W.01 contract.

The strongest contribution is the boundary, not the statistic. Because election
counts often violate Benford assumptions for innocent administrative reasons,
first-digit, second-digit, and last-digit tests should be treated as low-weight
screens. They may help decide where to look next. They do not establish what
happened, and they must never be presented as fraud evidence or certification
evidence on their own.
